\chapter{Obstructions, Equivalence, and Hallucination}

One of the repository's strongest unifying ideas is that several seemingly different
problems become gluing problems once the right local sections are chosen.

\section*{Bugs as obstruction classes}
In the sheaf-cohomological line of work, a bug is represented as a non-gluing local fact.
A local section at one site may allow a state that another site's viability condition
forbids. The disagreement lives on an overlap and can be collected into a $1$-cocycle.
When that cocycle is not a coboundary, it determines a nonzero class in first cohomology.
This is the precise sense in which the repository says ``bugs are $H^1$ classes.''

\begin{example}[Normalization bug]
Suppose an argument site allows an empty list while a downstream division site requires a
nonzero sum. The call and error sites disagree on their overlap. The mismatch is a
$1$-cocycle; if no local strengthening already present in the cover resolves it, the class
is nonzero and the bug is genuinely obstructive.
\end{example}

The slogan becomes especially valuable once it is quantified. The repository's older paper
material argues that the rank of $\Hc{1}$ over $\mathbb{F}_2$ corresponds to the minimum
number of independent fixes required to resolve the current family of gluing failures.
Whether or not every implementation path computes that exact number in the same way, the
architectural point is powerful: non-gluing is not just yes/no, it has structure and can
be decomposed.

\section*{Equivalence as gluing of local isomorphisms}
The equivalence machinery in \repo{src/deppy/equivalence/} makes the same move for program
comparison. Two programs are not compared by one magical global theorem from nowhere.
Instead, one constructs local correspondences, local semantic differences, and overlap
conditions. If the local isomorphisms glue, one obtains a global equivalence verdict. If
they do not, one obtains a localized reason for failure.

This reuse of the same geometry is one of the best reasons to take the repository's
mathematical self-description seriously. It would be much easier to implement bug finding
and equivalence with unrelated vocabularies. The decision not to do so shows that the
local-to-global picture is doing real architectural work.

\section*{Hallucination as type inhabitation failure}
Anti-hallucination is where the hybrid split between structural and semantic predicates
becomes especially visible. A fabricated API is a structural failure: the relevant symbol,
signature, or object does not exist. A semantic drift is a meaning failure: the code might
be structurally well-formed while still failing the intended rubric or behavior. A
contradiction between multiple checks is an inconsistency failure.

This is a very strong and very repository-specific claim. It says that hallucination should
not be treated as mystical stochastic misbehavior. It should be treated as typed failure.
The artifact has a required hybrid type. Structural checks attempt to inhabit the
structural component. Semantic checks attempt to inhabit the semantic component. Proof and
code layers contribute their own evidence. If the object fails to inhabit the required
space, the failure can be named precisely.

\begin{definition}[Hallucination as inhabitation failure]
Given an expected hybrid type $T$, a generated artifact $a$ hallucinates when $a$ fails to
inhabit $T$ at some required layer. Structural hallucinations fail structural inhabitation;
semantic hallucinations fail semantic inhabitation; contradictory hallucinations produce
incompatible evidence across layers.
\end{definition}

This is one of the clearest examples of the paper's title. Once local-to-global typed
structure is the common language, hallucination, equivalence, and bug finding are not
unrelated downstream products. They are different ways for a typed section to succeed or
fail to glue.
